\section{Beurteilung des Kursverlaufs}
\label{sec:kurs}

\subsection{Kurs und Auftragseingang}

\kurstabelle

Der Schlusskurs zum Jahresende 2025 betrug \je{\KursJahresende}{EUR}\Q{GeaGB}{14}; seither
ist er auf \je{\Kurs}{EUR} gestiegen (\KursDatum)\QW{Yahoo Finance, Kursreihe
G1A.DE}{quelle:kurse}{19.09.2026}.

\bk{Die beiden Ver\"anderungsspalten laufen nicht parallel. 2021 stieg der Kurs um ein
Vielfaches der \"Anderung des Auftragseingangs, 2022 fiel er, obwohl der Auftragseingang
zunahm, und 2024 stieg er stark bei fast unver\"andertem Auftragseingang. \"Uber den
gesamten Zeitraum sind beide Gr\"o{\ss}en gestiegen, von Jahr zu Jahr erkl\"art der
Auftragseingang den Kurs jedoch nicht. Das ist kein Mangel des Unternehmens: Der Kurs
enth\"alt, was der Markt \"uber die Jahre nach dem Berichtsjahr denkt, die Auftragszahl
nur das Berichtsjahr selbst.}

\subsection{Bewertung im Zeitverlauf}

\kgvtabelle

Das Kurs-Gewinn-Verh\"altnis betr\"agt heute \KgvHeute{} gegen einen Median von
\KgvMedian{} in den Jahren \KgvJahre; der heutige Wert liegt \"uber \Pz{\KgvRang} der
erfassten Jahresendwerte. Der freie Zufluss von \Mio{\FcfVoll}{EUR} entspricht
\FcfRendite{} des B\"orsenwerts, das Verh\"altnis von B\"orsenwert zu Buchwert des
Eigenkapitals betr\"agt \KBV.

\bk{Das KGV liegt \"uber dem eigenen Median, aber nicht au{\ss}erhalb der eigenen
Geschichte; die Jahre 2020 und 2021 waren teurer. Wer allein darauf schaut, sieht eine
normal bewertete Aktie. Die Zahl, die dieses Bild ver\"andert, ist der freie Zufluss: Er
verzinst den heutigen Kaufpreis mit \FcfRendite. Das Ergebnis je Aktie w\"achst seit 2020
in jedem Jahr, der freie Zufluss nicht, weil GEA in den vergangenen Jahren mehr investiert
hat, als das steigende Ergebnis an zus\"atzlichem Mittelzufluss brachte. Die Bewertung
folgt dem Ergebnis, die Verzinsung des Anlegers folgt dem Zufluss.}
